\documentclass[a4paper]{amsart}
\usepackage[french]{babel}
\usepackage[utf8]{inputenc}
\author{S. Puechmorel}
\date{\today}
\title{Déroulement du cours}
\begin{document}
\section{Généralités}
Le cours est divisé en neuf séances de deux heures. Tous les exercices, apparaissant dans le corps du texte sont à faire avec les
étudiants, à l'exception des numéros un et deux du chapitre "Topologie". Le format de l'évaluation est le suivant:
\begin{itemize}
    \item Six questions de cours couvrant l'ensemble des thèmes du programme, pour un total de six points. Elles seront extraites d'une liste
    distribuée aux étudiants deux semaines avant l'examen.
    \item Un exercice similaire à ceux vus durant le module, crédité de six points. 
    \item Un exercice original, noté sur huit points.    
\end{itemize}
\section{Topologie}
Seules les parties sur la compacité, la connexité simple et la connexité par arcs sont à traiter. Les exercices proposés sont destinés
au travail personnel des étudiants et ne font pas partie du programme de l'examen. Aucune démonstration n'est à faire en séance.
\section{holomorphie}
Les rappels de calcul différentiel ne sont pas à présenter, le cours proprement dit débute en section 2.2. 
\section{Integration complexe}
\section{Théorème des résidus}
\section{transformation de Laplace}
\end{document}